\documentclass[12pt]{ctexart}

\usepackage{amsmath,amssymb}
\usepackage{bm}
\usepackage{geometry}
\usepackage{hyperref}
\geometry{a4paper,margin=2.5cm}

\title{从三维重建到世界模型的技术路径综述\\[4pt]\large 几何、辐射场与显式表示}
\author{3D\_Reconstruction\_ing 项目笔记}
\date{\today}

\begin{document}

\maketitle

\begin{abstract}
本综述基于 \texttt{3D\_Reconstruction\_ing} 仓库中的文档与代码，系统梳理从传统三维重建到面向世界模型（World Model）的 3D/360 表示技术路径。
文章从几何层（几何一致性）、辐射场层（NeRF 及其变体）、显式高效表示层（3D Gaussian Splatting / 4DGS），以及语义与世界记忆接口四个层面展开，
给出关键数学公式与算法要点，为后续工程实践文档与实验提供理论背景。
\end{abstract}

\tableofcontents

\section{引言：三维重建与世界模型}

三维与全景重建（3D / 360 Reconstruction）关注从多视角图像或视频中恢复世界的几何结构与全景语义。
在世界模型视角下，我们期望的不只是离线重建一个静态场景，而是构建一个
\emph{可查询、可渲染、可推理} 的空间表示，用于支撑导航、编辑、理解与生成等多种下游任务。

本仓库的总体目标可以概括为：
\begin{quote}
围绕 Mip-NeRF 360 等真实世界数据，搭建从 COLMAP $\to$ Nerfstudio $\to$ 3DGS / 4DGS 的动手实践与技术验证管线，为“世界模型”的 3D / 360 表示打基础。
\end{quote}

从技术栈角度，核心可拆解为三层：
\begin{enumerate}
  \item \textbf{几何层（Geometry / SfM / MVS）：}\\
  估计相机位姿与稀疏/稠密点云，保证几何与物理一致性。
  \item \textbf{辐射场层（Radiance Field / NeRF）：}\\
  学习连续体积场函数 $f(\bm{x},\bm{d})$，从任意视角渲染高质量图像。
  \item \textbf{显式高效表示层（3DGS / 4DGS 等）：}\\
  用高斯云等显式结构替代纯 MLP 体渲染，在保证质量的前提下提升渲染效率与可编辑性。
\end{enumerate}

在此之上，再叠加语义层与统一的场景状态（Scene State），构成可查询、可组合的世界记忆。

\section{几何层：多视图几何与 COLMAP}

\subsection{针孔相机模型与坐标变换}

在经典多视图几何中，针孔相机模型将 3D 点 $\bm{X} = (X,Y,Z)^\top$ 映射到图像平面像素 $\bm{p} = (u,v)^\top$：
\begin{equation}
  \lambda
  \begin{bmatrix}
    u \\ v \\ 1
  \end{bmatrix}
  =
  \bm{K}
  \begin{bmatrix}
    \bm{R} & \bm{t} \\
    \bm{0}^\top & 1
  \end{bmatrix}
  \begin{bmatrix}
    X \\ Y \\ Z \\ 1
  \end{bmatrix},
  \quad
  \bm{K} =
  \begin{bmatrix}
    f_x & 0   & c_x \\
    0   & f_y & c_y \\
    0   & 0   & 1
  \end{bmatrix},
\end{equation}
其中 $\bm{R}\in SO(3)$、$\bm{t}\in\mathbb{R}^3$ 为相机外参，$\bm{K}$ 为内参，$\lambda$ 为尺度因子。
当考虑径向/切向畸变时，通常在像平面上先对归一化坐标做畸变矫正，再乘以内参。

\subsection{对极几何与本质矩阵}

在两视图情形下，空间点 $\bm{X}$ 在两个相机中的像素分别为 $\bm{p}_1,\bm{p}_2$（齐次坐标），
则存在本质矩阵 $\bm{E}$ 与基础矩阵 $\bm{F}$ 使得：
\begin{equation}
  \bm{p}_2^\top \bm{F} \bm{p}_1 = 0,\qquad
  \bm{F} = \bm{K}_2^{-\top} \bm{E} \bm{K}_1^{-1},\qquad
  \bm{E} = [\bm{t}]_\times \bm{R},
\end{equation}
其中 $[\bm{t}]_\times$ 为由平移向量 $\bm{t}$ 构造的反对称矩阵。
在 COLMAP 等 SfM 系统中，通常使用带 RANSAC 的 8 点或 5 点算法从匹配点对中估计 $\bm{F}/\bm{E}$，
再分解得到相对位姿 $(\bm{R},\bm{t})$（含尺度歧义）。

\subsection{三角化与 Bundle Adjustment}

已知相机投影矩阵 $\bm{P}_1,\bm{P}_2$ 与匹配像素 $\bm{p}_1,\bm{p}_2$，可通过线性 DLT 或中点法计算 3D 点
$\bm{X}$，这一过程称为三角化：
\begin{equation}
  \bm{p}_i \propto \bm{P}_i \bm{X}, \quad i=1,2.
\end{equation}
对整个场景而言，带噪声的观测会导致误差累积，需要通过捆绑调整（Bundle Adjustment, BA）进行全局非线性优化。
设 $\pi(\cdot)$ 为投影函数（含内参和畸变），$\bm{p}_{ij}$ 为第 $i$ 台相机观测到的第 $j$ 个点的像素位置，则
BA 的典型目标函数为：
\begin{equation}
  \min_{\{\bm{R}_i,\bm{t}_i\},\{\bm{X}_j\}}
  \sum_{i,j} \rho\Bigl(
    \bigl\|
      \pi(\bm{R}_i \bm{X}_j + \bm{t}_i) - \bm{p}_{ij}
    \bigr\|^2
  \Bigr),
\end{equation}
其中 $\rho(\cdot)$ 为鲁棒核函数（如 Huber），以降低误匹配对优化的影响。
COLMAP 采用增量式 SfM：从少量视图初始化，逐步加入新视图并在每一步执行局部或全局 BA。

\subsection{多视图立体（MVS）与稠密重建}

在 SfM 获得准确相机位姿后，多视图立体（MVS）进一步估计每个像素的深度，从而形成稠密点云。
COLMAP 中典型流程为：
\begin{enumerate}
  \item 去畸变与统一成像模型：将图像统一到理想针孔模型，方便后续几何推理；
  \item PatchMatch Stereo：为每个像素估计局部平面（由深度 $d$ 与法向 $\bm{n}$ 描述），
        通过随机初始化、邻域传播与随机扰动迭代优化，
        在光度一致性损失（如 NCC 或绝对差）下收敛；
  \item 多视图深度融合：对各图像的深度图进行几何一致性检测，
        将多视图一致的像素融合为 3D 点，得到稠密点云。
\end{enumerate}

在本仓库面向 NeRF / Nerfstudio 的实践中，\emph{稀疏层的相机位姿已经足够驱动后续辐射场训练}，
稠密 MVS 是可选的增强步骤。

\section{辐射场层：NeRF 与体渲染}

\subsection{连续辐射场表示}

NeRF 将场景表示为一个连续的辐射场函数
\begin{equation}
  f : (\bm{x}, \bm{d}) \mapsto (\bm{c}, \sigma),
\end{equation}
其中 $\bm{x}\in\mathbb{R}^3$ 为 3D 空间点，$\bm{d}\in\mathbb{S}^2$ 为视线方向，
$\bm{c}\in\mathbb{R}^3$ 为 RGB 颜色，$\sigma\in\mathbb{R}_{\ge 0}$ 为体密度。
在原始 NeRF 中，$f$ 由一个多层感知机（MLP）拟合，输入为对 $\bm{x},\bm{d}$ 做位置编码（positional encoding）后的高维特征：
\begin{equation}
  \gamma_L(x) = \Bigl(
    \sin(2^0\pi x), \cos(2^0\pi x),\dots,
    \sin(2^{L-1}\pi x), \cos(2^{L-1}\pi x)
  \Bigr),
\end{equation}
对每个坐标维度独立编码后再拼接。
本仓库的语义 NeRF 模型也采用了类似的正余弦位置编码结构。

\subsection{体渲染积分与离散近似}

给定一条射线
\begin{equation}
  \bm{r}(t) = \bm{o} + t \bm{d}, \quad t \in [t_n, t_f],
\end{equation}
其中 $\bm{o}$ 为相机中心，$\bm{d}$ 为单位方向向量，
NeRF 定义该射线对应像素颜色的体渲染积分为：
\begin{equation}
  C(\bm{r}) = \int_{t_n}^{t_f}
    T(t)\,\sigma\bigl(\bm{r}(t)\bigr)\,
    \bm{c}\bigl(\bm{r}(t),\bm{d}\bigr)\,\mathrm{d}t,
  \quad
  T(t) = \exp\left(
    -\int_{t_n}^{t} \sigma(\bm{r}(s))\,\mathrm{d}s
  \right).
\end{equation}

实际实现中，将积分离散化为若干采样点 $\{t_i\}$，得到近似：
\begin{equation}
  \hat{C}(\bm{r}) =
  \sum_i T_i \bigl(1 - \exp(-\sigma_i \delta_i)\bigr)\bm{c}_i,\quad
  T_i = \exp\Bigl(-\sum_{j<i} \sigma_j \delta_j\Bigr),
\end{equation}
其中 $\delta_i = t_{i+1}-t_i$ 为相邻采样点间距。
这一形式既保持了可微性，也较好逼近了体积渲染物理过程。

\subsection{训练目标与变体}

在监督训练设置下，给定一组多视角图像与相机位姿，
对每条射线 $\bm{r}$，计算预测颜色 $\hat{C}(\bm{r})$，与真实像素颜色 $\bm{C}^\ast(\bm{r})$ 做损失：
\begin{equation}
  \mathcal{L}_\text{RGB}
  =
  \frac{1}{N}\sum_{\bm{r}}
    \bigl\| \hat{C}(\bm{r}) - \bm{C}^\ast(\bm{r}) \bigr\|_2^2
  \quad \text{或更一般的感知损失}.
\end{equation}

Nerfstudio 中的 \texttt{nerfacto} 等变体在以下方面做了大量工程优化：
\begin{itemize}
  \item 使用多分辨率哈希网格等高效编码替代纯 MLP；
  \item 采用分层/重要性采样，提高采样点利用率；
  \item 加入外观编码、曝光校正等，以增强真实感与泛化能力。
\end{itemize}

尽管具体结构各异，它们在数学上仍服从 ``连续场 + 体渲染 + 可微优化'' 这一统一范式。

\section{显式高效表示：3D Gaussian Splatting 与 4D 扩展}

\subsection{3D 高斯表示与参数化}

3D Gaussian Splatting（3DGS）使用一组带属性的三维高斯来显式表示场景。
对单个高斯，可写为
\begin{equation}
  \mathcal{G}(\bm{x})
  = \alpha \cdot \exp\!\Bigl(
      -\tfrac{1}{2}
      (\bm{x}-\bm{\mu})^\top
      \bm{\Sigma}^{-1}
      (\bm{x}-\bm{\mu})
    \Bigr),
\end{equation}
其中 $\bm{\mu}\in\mathbb{R}^3$ 为中心，$\bm{\Sigma}\in\mathbb{R}^{3\times 3}$ 为正定协方差矩阵，
$\alpha$ 为不透明度或密度相关系数，颜色可由 RGB 或球谐系数参数化。

为了保证 $\bm{\Sigma}$ 正定，实践中常采用
\begin{equation}
  \bm{\Sigma} = \bm{R}\bm{S}\bm{S}^\top \bm{R}^\top,
\end{equation}
其中 $\bm{R}\in SO(3)$ 为旋转矩阵，$\bm{S}=\mathrm{diag}(s_x,s_y,s_z)$ 为尺度矩阵。
3DGS 的优化过程即是在多视角图像监督下对
$(\bm{\mu},\bm{R},\bm{S},\alpha,\text{color})$ 等参数做梯度更新。

\subsection{Splatting 与可微渲染}

在渲染时，首先将三维高斯按当前相机位姿投影到二维像素平面，
得到 2D 高斯（中心与协方差经过线性变换），再按照深度从远至近排序，
对每个像素做 alpha blending：
\begin{equation}
  \bm{C}_\text{out} =
  \sum_{k}
    \tilde{\alpha}_k
    \Bigl(
      \prod_{j<k} (1-\tilde{\alpha}_j)
    \Bigr)
    \bm{c}_k,
\end{equation}
其中 $\tilde{\alpha}_k$ 为第 $k$ 个高斯在该像素处的透明度权重，
$\bm{c}_k$ 为其颜色。
这一过程是可微的，可直接用反向传播优化高斯参数。

与 NeRF 相比，3DGS 在保持较高重建质量的同时，
大幅提高了渲染速度与交互效率，更适合实时查看、编辑与下游任务。

\subsection{4DGS 与动态场景}

4D Gaussian Splatting（4DGS）在 3DGS 的基础上将时间或隐空间纳入参数，
例如令高斯中心或颜色为时间的函数：
\begin{equation}
  \bm{\mu}(t),\quad
  \bm{\Sigma}(t),\quad
  \bm{c}(t),
\end{equation}
从而表示动态场景或动作序列。
在世界模型语境下，4DGS 扮演了 ``状态随时间演化'' 的显式表征角色，
有利于对时空一致性的建模与推理。

\section{语义层与世界记忆接口}

\subsection{2D 语义分割与像素级监督}

语义层的第一步通常是对每张输入图像做 2D 语义分割。
设某张图像尺寸为 $H\times W$，类别数为 $K$，
语义分割网络输出每个像素的 logits：
\begin{equation}
  \bm{z}_{uv} \in \mathbb{R}^K,\quad (u,v)\in\{1,\dots,W\}\times\{1,\dots,H\}.
\end{equation}
经过 softmax 得到类别概率
\begin{equation}
  p(y_{uv}=k) = \frac{\exp(z_{uv}^{(k)})}{\sum_{k'} \exp(z_{uv}^{(k')})},
\end{equation}
训练时使用交叉熵损失：
\begin{equation}
  \mathcal{L}_\text{seg}
  = - \frac{1}{HW}
    \sum_{u,v} \log p\bigl(y_{uv}^\ast\bigr),
\end{equation}
其中 $y_{uv}^\ast$ 为标注的真值类别。

本仓库的 \texttt{scripts/run\_semantic\_labeling.py} 负责对 Nerfstudio 预处理后的图像
生成像素级 label map，并可选汇总为每张图的类别统计等元数据。

\subsection{3D 语义融合：从图像到点云}

在已知 COLMAP 稀疏点云与其多视图观测的前提下，
可将 2D 语义标签投票到 3D 点上，得到 ``带语义的点云''。
设某个 3D 点 $\bm{X}$ 在若干图像中的投影像素为 $\{(u_i,v_i)\}$，
对应的 2D 语义标签为 $\{y_i\}$，则可定义该点的 3D 语义标签为
\begin{equation}
  \hat{y}(\bm{X})
  = \arg\max_{k}
    \sum_{i} w_i \,\mathbf{1}[y_i = k],
\end{equation}
其中 $\{w_i\}$ 为基于视角、深度或置信度的权重。

在本仓库中，\texttt{semantic\_scene.json} 记录了这样的 3D 语义点信息，
并提供了按类别、按 3D 区域、按查询点检索的接口。

\subsection{语义 NeRF 与语义 3DGS}

在此基础上，可以在辐射场与 3DGS 层面进一步引入语义监督。

\paragraph{语义 NeRF：}
对每个采样点，NeRF 模型除了输出 $(\bm{c},\sigma)$ 外，还预测语义 logits $\bm{s}\in\mathbb{R}^K$。
在体渲染过程中，同时对 RGB 与语义进行加权积分：
\begin{align}
  \hat{\bm{C}}(\bm{r}) &= \sum_i w_i \bm{c}_i, \\
  \hat{\bm{s}}(\bm{r}) &= \sum_i w_i \bm{s}_i,
\end{align}
其中 $w_i$ 为与密度相关的体渲染权重。
对每条射线的语义标签 $y^\ast$ 使用交叉熵损失：
\begin{equation}
  \mathcal{L}_\text{sem-NeRF}
  =
  - \frac{1}{N}
    \sum_{\bm{r}}
    \log
    \frac{\exp\bigl(\hat{s}_{y^\ast}(\bm{r})\bigr)}
         {\sum_k \exp\bigl(\hat{s}_k(\bm{r})\bigr)}.
\end{equation}
总损失为 RGB 与语义损失的加权和。
本仓库的 \texttt{src/semantic\_nerf} 模块实现了一个轻量级语义 NeRF。

\paragraph{语义 3DGS：}
类似地，可在每个高斯上附加语义 logits，渲染语义图并用 2D 标签监督。
由于几何与颜色已由基础 3DGS 学得，语义头的训练可以在冻结几何与 RGB 的条件下进行，
从而提高优化稳定性。

\section{统一静态场景状态与世界模型视角}

为了在世界模型实验中复用同一真实场景，本仓库通过
\emph{静态场景状态（Scene State）} 将多种表示统一到一个 JSON 描述中，
例如 \texttt{mipnerf360/db/playroom/world\_state.playroom.json}。
其核心思想包括：
\begin{itemize}
  \item 统一的场景 ID、根目录与世界坐标系说明；
  \item 集中记录 \texttt{COLMAP}、\texttt{Nerfstudio/NeRF}、
        \texttt{3DGS} 与 \texttt{semantics} 各自的数据与模型路径；
  \item 通过 \texttt{SceneState} 与 \texttt{SemanticScene} 等数据结构，
        提供按类别、按 3D 区域、按点的可查询接口。
\end{itemize}

从更宏观的角度看，\emph{几何层保证物理一致性，辐射场层保证真实感，
显式层保证效率与可编辑性，语义层与 Scene State 则将这些表示组织成统一的世界记忆单元}。
这正是面向世界模型的 3D/360 重建技术路径的核心理念。

\section{结语}

本文以 \texttt{3D\_Reconstruction\_ing} 仓库为背景，围绕几何、辐射场、显式高斯表示与语义层，
梳理了从传统三维重建走向世界模型的关键数学原理与算法框架。
配合工程实践文档（另见配套的工程实现教程 \LaTeX 文档），
读者可以一方面从理论上理解各层组件的角色与联系，
另一方面在真实数据集（如 Mip-NeRF 360）上动手打通整条 Pipeline。

\end{document}

